\chapter{Compilers as Incremental Databases}
\label{ch:incremental-databases}

Changing a function body from returning $7$ to returning $8$ requires new executable code, but may leave its type, calling convention, and callers' type checks unchanged.
A compiler that treats the edit as a command to rerun every pass loses this distinction.
A compiler that caches only by filename loses it too: the file changed, but many answers about it did not.
The useful unit of reuse is a question with explicit inputs and a result whose meaning is stable across revisions.

The database perspective is an organization of compiler state, not a requirement to store every instruction in SQL.
It asks which facts are authoritative, which are derived, how dependencies are recorded, and which revision a result describes.
SSA already exposes value dependencies; the next step is to make dependencies between compiler computations equally explicit.

\section{Separate facts, queries, and snapshots}
\label{sec:ch61-schema}

A \term{compiler query} is a deterministic computation identified by a query kind and a key, such as the type of a stable function identity.
Its provider reads inputs and other queries through a tracked interface.
The provider must not secretly consult a mutable global option, the current directory, or a wall clock.
Such a read would become an invisible dependency and could make a cached answer incorrect.
Target features, language modes, imported interfaces, and compiler configuration therefore belong in the input model.

Table~\ref{tab:ch61-schema} sketches a relational view of a small IR.
Operand position matters: subtracting operand zero from operand one is not the same computation as the reverse.
Block membership and control successors remain explicit, even when instruction nodes are shared or interned.
A database schema does not remove dominance, typing, or effect-order invariants.
Those are integrity conditions checked over a snapshot.

\begin{table}[htbp]
\caption{A relational IR separates stable identities from ordered references.}
\label{tab:ch61-schema}
\centering
\begin{tabular}{ll}
\toprule
Relation & Representative fields \\
\midrule
Instruction & identity, opcode, type, owning block \\
Operand & user identity, position, definition identity \\
Successor & source block, edge position, target block \\
Origin & IR identity, source identity, transformation identity \\
Query result & key, value, dependencies, revision stamps \\
\bottomrule
\end{tabular}
\end{table}

A \term{snapshot} fixes the authoritative inputs seen by one evaluation.
An edit creates a new revision instead of mutating inputs halfway through a provider.
Stable identities are not byte offsets: inserting a comment should not necessarily rename every function below it.
Conversely, preserving an identity must not falsely claim that its contents stayed unchanged.
Identity and content fingerprints solve different problems.

\section{Validate dependencies before recomputing}
\label{sec:ch61-validation}

The companion stores two stamps for each result: the revision in which its value last changed, and the revision in which it was last verified.
Suppose a caller query was verified at revision zero.
At revision one its type dependency is recomputed but still returns \code{int}.
The type result's verification stamp advances while its change stamp stays zero.
The caller can therefore reuse its old result.
This is \term{early cutoff}: an input change stops propagating when a derived value is unchanged.
The Rust compiler guide describes a production red-green validation scheme with this central distinction~\cite{rust-2026-incremental}.

Algorithm~\ref{alg:ch61-query} assumes pure providers, immutable result values, exact equality, and an acyclic query graph within each evaluation.
The dependency list is ordered by the provider's first read of each key.
Validation stops at its first changed dependency.
Continuing through the rest of the old list could evaluate a query that the new control path never requests and that is no longer valid.

\begin{algorithm}[htbp]
\caption{Demand a query in the current immutable revision.}
\label{alg:ch61-query}
\begin{pseudocode}
\Require key $q$, revision $r$, pure tracked providers, cache
\Ensure result equal to fresh evaluation in revision $r$
\State record $q$ as a dependency of its caller
\If{$q$ is an input}
  \State \Return current input value
\EndIf
\State reject a recursive request for an active key
\If{$q$ was verified in $r$}
  \State \Return cached value
\EndIf
\State mark $q$ active
\If{an old result exists}
  \State validate old dependencies in read order, stopping at the first changed value
  \If{all dependencies are unchanged since the old verification}
    \State set verification stamp to $r$; unmark $q$; \Return old value
  \EndIf
\EndIf
\State run the provider while recording a fresh dependency list
\State retain the old change stamp if the result is equal; otherwise use $r$
\State atomically replace the cached result and dependencies; unmark $q$
\State \Return new value
\end{pseudocode}
\end{algorithm}

For $V$ reached keys and $E$ reached dependency edges, one revision's validation traversal costs $O(V+E)$ plus provider execution and value-comparison costs, assuming expected constant-time map access.
Each key is verified once per revision; a request for an already verified key is a cache lookup.
This bound does not make providers cheap: parsing a changed file or comparing a large exact result still costs time.
The teaching implementation uses exact Python equality rather than pretending that a fixed-width hash is a collision-free proof.
A production implementation needs a documented fingerprint policy and stable serialization.

Figure~\ref{fig:ch61-cutoff} shows why the changed body does not force the client to run.
Table~\ref{tab:ch61-revisions} gives the actual values and cumulative provider counts.
The third edit changes the type and does reach the client.

\begin{figure}[htbp]
\centering
\begin{tikzpicture}[node distance=16mm]
\node[basic block] (body) {body: 7 becomes 8};
\node[basic block,below=of body] (kind) {kind: int remains int};
\node[basic block,below=of kind] (client) {client: reuse int-call};
\draw[cfg edge] (body)--node[right,font=\small] {recompute}(kind);
\draw[cfg edge,dashed] (kind)--node[right,font=\small] {unchanged answer}(client);
\end{tikzpicture}
\caption{Validation propagates a changed answer rather than merely a changed source file.}
\label{fig:ch61-cutoff}
\end{figure}

\begin{table}[htbp]
\caption{Exact result comparison suppresses one unnecessary client evaluation.}
\label{tab:ch61-revisions}
\centering
\begin{tabular}{llll}
\toprule
Body & Kind & Kind evaluations & Client evaluations \\
\midrule
$7$ & int & $1$ & $1$ \\
$8$ & int & $2$ & $1$ \\
text value & text & $3$ & $2$ \\
\bottomrule
\end{tabular}
\end{table}

\section{Queries can be reactive without being eager}
\label{sec:ch61-reactive}

An editor can subscribe to diagnostics while executable-code queries remain undemanded.
A runtime can request machine code only when an entry point becomes hot.
Both use the same invalidation mechanism but different roots of demand.
Calling this arrangement reactive does not require immediately rebuilding every invalidated artifact.
Invalidation records doubt; demand determines when to discharge it.

Cycles need an explicit policy.
Mutually recursive functions do not force type checking to be a cyclic cache recursion: declarations or strongly connected components can form the provider boundary.
A monotone analysis over a finite domain can instead expose a fixed-point provider with its own convergence contract.
Simply returning a partially computed cached value on recursion would silently change the analysis.
Provider failures must also unwind active-key markers, so one diagnostic does not poison later requests.

\section{Publish transformations as transactions}
\label{sec:ch61-transactions}

A relational optimizer may express a candidate as a join over opcode, operand, type, and effect facts.
Matching an addition of zero is easy; establishing that its replacement dominates all uses and preserves semantic flags is the harder part.
The query identifies candidates, while a checked transaction performs mutation.
Algorithm~\ref{alg:ch61-transaction} prevents a worker from committing a result computed from a stale revision.
Its validation work is proportional to the facts examined by the invariant checker, not necessarily to the number of changed rows.

\begin{algorithm}[htbp]
\caption{Commit an IR rewrite against the snapshot it actually inspected.}
\label{alg:ch61-transaction}
\begin{pseudocode}
\Require immutable snapshot $S$, proposed edits $D$, current revision owner
\Ensure either one valid new snapshot or no published change
\State construct candidate snapshot $S'$ from $S$ and $D$
\State check types, dominance, edge arity, and required effect ordering in $S'$
\State enter the serialized publication operation
\If{the current snapshot is not $S$}
  \State \Return stale result
\EndIf
\State publish $S'$ with new input stamps and recorded provenance
\State \Return committed
\end{pseudocode}
\end{algorithm}

Immutable snapshots also let analysis workers run concurrently without observing half-rewritten use lists.
An implementation may use persistent data structures or private copies; neither makes conflict resolution automatic.
Two individually legal transformations can disagree about the same definition.
Revalidation, merging with a proof, or retrying against the newer snapshot are distinct engineering choices.

\section{Provenance is not an inverse compiler}
\label{sec:ch61-provenance}

\term{Provenance} records where an artifact came from and which decisions produced it.
Inlining creates many machine locations from one source expression; common-subexpression elimination creates one value from several expressions.
Constant folding maps many distinct source trees to the same constant.
These many-to-one operations destroy the uniqueness needed for an inverse.
A \term{bidirectional transformation} therefore needs an explicit update policy and retained information, not merely arrows drawn both ways.

For example, folding both \code{2+3} and \code{1+4} gives $5$.
An edit of that result to $6$ does not specify which operand to change or whether to replace the whole expression.
Figure~\ref{fig:ch61-provenance} records both origins without choosing an invented source edit.
Debugger locations, explanation tools, and deoptimization maps each select different information from this relation.
Debug provenance can be approximate; reconstructing a running frame requires an executable state map, developed in Chapter~\ref{ch:speculative-state-transfer}.

A restricted bidirectional editor can define a view function $view(s)$ and an update function $put(s,v)$ that uses both the old source $s$ and a requested view $v$.
Two useful consistency requirements are $put(s,view(s))=s$ and $view(put(s,v))=v$ for admitted updates.
The first preserves an unchanged view without rewriting the source; the second requires an accepted update to be reflected in the next view.
Neither condition chooses a unique source edit when folding has erased structure.
The tool must specify which updates it admits and which source fragments it may replace.

For a constant-only view, replacing the whole expression by the requested literal satisfies the second requirement, but violates the first if applied unconditionally to an unchanged view of \code{2+3}.
A policy can preserve the old source when the value is unchanged and replace the expression otherwise.
That policy is deterministic and easy to describe, yet loses expression structure on edits.
Keeping richer provenance can support less destructive policies, but their behavior remains a tool contract rather than a recovered inverse of optimization.

\begin{figure}[htbp]
\centering
\begin{tikzpicture}[node distance=15mm and 18mm]
\node[basic block] (a) {source: 2+3};
\node[basic block,right=of a] (b) {source: 1+4};
\node[basic block,below=of a] (c) {IR constant 5};
\draw[cfg edge] (a)--(c); \draw[cfg edge] (b)--(c);
\node[basic block,below=of c] (p) {origin set: both expressions};
\draw[cfg edge] (c)--(p);
\end{tikzpicture}
\caption{Retaining multiple origins supports explanation without defining a unique inverse edit.}
\label{fig:ch61-provenance}
\end{figure}

\section{Implementation and exercises}
\label{sec:ch61-exercises}

The \code{Queries} class in \path{example-code/architecture-labs/labs.py} implements ordered validation, exact early cutoff, dependency replacement, and cycle rejection.
Tests compare incremental answers with fresh evaluators over $100$ input combinations.
They also remove a no-longer-reachable input to detect accidental evaluation of obsolete dependencies.
The model is single-threaded, keeps all results in memory, and requires inputs and results to remain immutable between edits.

\begin{enumerate}
\item\label{ex:ch61-cutoff} A query returns the parity of input $x$, and its client returns that parity plus $y$.
Starting with $(x,y)=(2,10)$, change $x$ to $4$, then $y$ to $11$.
Compute the client results and identify which providers rerun.
\item\label{ex:ch61-branch} A query reads a flag before conditionally reading a function body.
Explain why validating the old body after observing a changed flag can be incorrect.
\item\label{ex:ch61-hidden} A code-generation provider reads the target CPU from an untracked environment variable.
Construct an edit sequence that makes its cache unsound and repair the dependency model.
\item\label{ex:ch61-inverse} Give two source expressions that fold to the same constant and two different valid source edits for changing that constant.
State what additional policy a source-update tool needs.
\end{enumerate}

\section{Further reading}
\label{sec:ch61-reading}

The Rust compiler guide explains ordered dependency validation and why unchanged results cut off propagation~\cite{rust-2026-incremental}.
Firm illustrates how a graph representation exposes value and control relationships without requiring a relational storage engine~\cite{braun-2011-firm}.
Chapter~\ref{ch:continuous-compilation} extends snapshot ownership to asynchronous runtime compilation and artifact retirement.
